\section{Control Channel Fit}
\label{sec:ControlChannel}

The extraction of the branching fractions of the \Bmumu decay is performed with respect to a well-known reference channel, the \BpmKpmJpsi decay. However, there are few aspects related to the kinematic of the \Bds mesons which are dependent from the electric charge of the meson (such as, for instance, the $B$-candidate isolation and the fragmentation process). Therefore, an additional control channel is needed to study more closely the neutral $B$-meson kinematic and to verify that the corrections and the modelling of the \BpmKpmJpsi decay is suitable also for that purpose.\\
To reach this goal, the \BsJpsiPhi\ decay channel is used. The focus will be only on the \Bs meson, which is dominant with respect to the \Bd meson in the di-muon decay channel. 
\BsJpsiPhi\ events are selected as outlined in Section \ref{}. The main difference with respect to the reference channel is the requirement to have two kaons in the final state whose invariant mass is compatible with that of the $\phi$ meson. Both kaons are required to have $p_T$ $>$ 1 GeV, and their invariant mass is required to be interval $1005 \ \text{MeV} < m_\phi < 1035 \ \text{MeV}$. \\
An unbinned extended ML fit to the $J/\psi \phi$ invariant mass distribution in data is made to extract the \BsJpsiPhi\ signal component. The total PDF used in the fit is made of a superposition of two Gaussian to model the signal component and a third order Chebyshev polynomial to model the combinatorial background component. The dataset was collected using the same trigger strategy of the main analysis, collecting the same integrated luminosity as the \BpmKpmJpsi channel. Figure \ref{fig:jpsiphi_mass} shows the $J/\psi \phi$ invariant mass fit to the dataset used. \\
By means of the \textit{sPlot} tool, it was possible to extract the distribution of \BsJpsiPhi\ events in all 15 variables entering the continuum BDT discriminant. These distributions are compared with that extracted from simulation to verify that the \BsJpsiPhi\ MC sample well describe the data. All variables but the isolation $B_{iso}$ are well described in simulation...... Figure \ref{fig:isol_jpsiphi} shows the superposition of $B_{iso}$ distribution in data and MC for \BsJpsiPhi\ events. The ratio between this two distribution was used to reweight the \Bsmumu\ MC samples. The same plot for the remaining variables entering in the continuum BDT discriminant are shown in Appendix \ref{app:jpsiphi}.

\textred{PUT PLOT IMAGE AND SOME DATA MC COMPARISON}

\subsection{DDW corrections check}
 
Similarly to what has been done in Section \ref{DDW} to extract DDW corrections from data using \BpmKpmJpsi\ events, the same procedure is used in  \BsJpsiPhi\ events to extract the same corrections. The goal of this check is to prove that DDW corrections are independent from the $B$-meson used to evaluate them. Data / MC comparisons of the $p_T$ and $\eta$ distributions for \BsJpsiPhi\ events are shown in Figure \ref{fig:Jpsiphi_pt_eta} and then used to compute DDW corrections shown in Figure \ref{fig:DDWJpsiphi}. Finally, figure \ref{fig:DDWcomp} shows the superposition of the DDW corrections in $p_T$ and $\eta$ extracted in the two channels (i.e. \BpmKpmJpsi and \BsJpsiPhi) respectively. It is possible to notice that....(hopefully they'll be compatible!).
